\documentclass[../../../include/open-logic-section]{subfiles}

\begin{document}

\olfileid{sth}{choice}{vitali}
\olsection{Lampiran: Paradoks Vitali}

Supaya bisa netepaké apa konstruksi Banach--Tarski bisa ditampa,
kita perlu nliti \emph{buktiné}. Sayangé, iki mbutuhaké aljabar
kang luwih akèh tinimbang kang bisa kita aturaké ing kéné.
Nanging kita bisa mènèhi sawatara cathetan cekak kang nerangaké
inti bukti Banach--Tarski,\footnote{Kanggo katrangan luwih jembar,
delengen \cite{Weston2003} utawa \cite{Wagon2016}.} kanthi
nggatekaké asil iki:

\begin{thm}[Paradoks Vitali (ing $\ZFC$)]\ollabel{vitaliparadox}
Sembarang bunder bisa dipérang dadi pérangan kaétung cacahé;
pérangan-pérangan iku bisa disusun maneh (kanthi rotasi lan translasi)
dadi rong salinan bunder mau.
\end{thm}

Paradoks Vitali luwih gampang dibuktèkaké tinimbang Paradoks
Banach--Tarski. Kita nyebut ``Paradoks Vitali'' amarga asil iki
katut saka konstruksi himpunan kang ora bisa diukur déning Vitali
ing \citeyear{Vitali1905}. Nanging pérangan teori himpunan ing
kaloro buktiné padha banget. Béda wigatiné yaiku Banach--Tarski
mbagi dadi pérangan \emph{winates} cacahé, déné Paradoks Vitali
mbagi dadi pérangan \emph{tanpa wates nanging kaétung}. Miturut
\citet{Weston2003}, Paradoks Vitali ``ora nganti semono
nggumunaké kaya Paradoks Banach--Tarski, nanging nuduhaké yèn
paradoks géometris bisa muncul sanadyan ing kahanan kang prasaja.''

Paradoks Vitali ngrembug wujud rong dimensi, yaiku bunder.
Mula kita makarya ing bidang $\Real^2$. Anggepen $\rotationsgroup$
iku himpunan rotasi titik-titik ngubengi titik asal kanthi arah
jarum jam, kanthi sudut kang kelipatan \emph{rasional} saka puteran wutuh
lan diwakili ing $[0,2\pi)$. Ing ngisor iki ana sawatara kasunyatan
aljabar babagan $\rotationsgroup$; bukti bakal njlentrehaké tegesé:

\begin{lem}\ollabel{rotationsgroupabelian}
$\rotationsgroup$ mbentuk {grup} abelian miturut komposisi fungsi.
\end{lem}

\begin{proof}
Tulis $0_{\rotationsgroup}$ kanggo rotasi kanthi sudut $0$ radian.
Iki unsur identitas ing $\rotationsgroup$, amarga
$\comp{0_{\rotationsgroup}}{\rho} = \comp{\rho}{0_{\rotationsgroup}} =
\rho$ kanggo sembarang $\rho \in \rotationsgroup$.

Saben unsur nduwèni invers. Rotasi identitas dadi inversé dhéwé.
Yèn $\rho \in \rotationsgroup$ muter kanthi sudut $r$ radian kang
ora nol, mula $\rho^{-1} \in \rotationsgroup$ muter kanthi sudut
$2\pi - r$ radian, saéngga $\rho \circ \rho^{-1} = 0_\rotationsgroup$.

Komposisi asosiatif: $\comp{\rho}{(\comp{\sigma}{\tau})} =
\comp{(\comp{\rho}{\sigma})}{\tau}$ kanggo sembarang
$\rho, \sigma, \tau \in \rotationsgroup$.

Komposisi komutatif: $\comp{\rho}{\sigma} =
\comp{\sigma}{\rho}$ kanggo sembarang $\rho, \sigma \in
\rotationsgroup$.
\end{proof}

Pancen, grup $\rotationsgroup$ bisa dipérang dadi rong pérangan,
lan salah siji pérangan waé cukup kanggo mbalèkaké grup wutuh:

\begin{lem}\ollabel{disjointgroup}
Ana partisi $\rotationsgroup$ dadi rong himpunan kang pisah,
$\rotationsgroup_{1}$ lan $\rotationsgroup_{2}$, kang saben-saben
dadi basis kanggo $\rotationsgroup$.
\end{lem}

\begin{proof}
Anggepen $\rotationsgroup_{1}$ ngemot rotasi kanthi sudut
kelipatan rasional saka puteran wutuh ing $[0, \pi)$; tetepna
$\rotationsgroup_{2} = \rotationsgroup \setminus \rotationsgroup_1$.
Miturut aljabar dhasar,
$\Setabs{\comp{\rho}{\rho}}{\rho \in \rotationsgroup_1} =
\rotationsgroup$. Asil kang padha uga lumaku kanggo
$\rotationsgroup_2$.
\end{proof}

Kasunyatan grup iki bakal kita gunakaké kanggo mbuktekaké
\olref{vitaliparadox}. Anggepen $\onesphere$ bunder satuan,
yaiku himpunan titik kang jaraké pas~$1$ saka titik asal bidang:
$\Setabs{\tuple{r,s} \in \Real^2}{\sqrt{r^2+s^2}=1}$.
Kita bakal mbagi $\onesphere$ kanthi relasi iki ing $\onesphere$:
\[
	r \sim s \emph{ yèn lan mung yèn }(\exists \rho \in \rotationsgroup)\rho(r) = s.
\]
Tegesé, rong titik ing $\onesphere$ gegandhengan yèn siji titik
bisa dipindhah menyang sijiné kanthi rotasi sudut rasional kali
puteran wutuh ngubengi titik asal. Ora nggumunaké yèn:

\begin{lem}
$\sim$ iku relasi ekuivalensi.
\end{lem}

\begin{proof}
Cetha saka \olref{rotationsgroupabelian}.
\end{proof}

Saiki kita migunakaké Pilihan kanggo ngasilaké himpunan $C$ kang
ngemot pas siji unsur saka saben kelas ekuivalensi $\onesphere$
miturut $\sim$. Tegesé, kita njupuk fungsi pilihan $f$ ing himpunan
kelas ekuivalensi,\footnote{Amarga $\rotationsgroup$ !!{enumerable},
saben !!{element} saka $E$ uga !!{enumerable}. Amarga $\onesphere$
!!{nonenumerable}, saka \olref{vitalicover} lan
\olref[card-arithmetic][simp]{kappaunionkappasize} kita ngerti yèn
$E$ !!{nonenumerable}. Mula ing kéné digunakaké Pilihan kang
\emph{ora} kaétung.}
\[
	E = \Setabs{\equivrep{r}{\sim}}{r \in \onesphere},
\]
lan tetepna $C = \ran{f}$. Kanggo saben rotasi
$\rho \in \rotationsgroup$, himpunan $\funimage{\rho}{C}$
ngemot titik-titik asil nrapaké $\rho$ marang saben titik ing $C$.
Rong asil sabanjuré nuduhaké yèn himpunan-himpunan iki nutupi
bunder kanthi jangkep lan tanpa tumpang tindih:

\begin{lem}\ollabel{vitalicover}
$\onesphere = \bigcup_{\rho \in \rotationsgroup} \funimage{\rho}{C}$.
\end{lem}

\begin{proof}
Tetepna $s \in \onesphere$. Ana $r \in C$ kang nyukupi
$r \in \equivrep{s}{\sim}$, yaiku $r \sim s$, saéngga
$\rho(r) = s$ kanggo sawijining $\rho \in \rotationsgroup$.
\end{proof}

\begin{lem}\ollabel{vitalinooverlap}
Yèn $\rho_1 \neq \rho_2$, mula
$\funimage{\rho_1}{C} \cap \funimage{\rho_2}{C} = \emptyset$.
\end{lem}

\begin{proof}
Anggepen $s \in \funimage{\rho_{1}}{C} \cap
\funimage{\rho_{2}}{C}$. Mula
$s = \rho_{1}(r_{1}) = \rho_{2}(r_{2})$ kanggo sawijining
$r_{1}, r_{2} \in C$. Saka kéné
$\rho^{-1}_2(\rho_1(r_1)) = r_2$, lan
$\comp{\rho_1}{\rho^{-1}_2} \in \rotationsgroup$, saéngga
$r_{1} \sim r_{2}$. Mula $r_1 = r_2$, amarga $C$ milih pas siji
unsur saka saben kelas ekuivalensi miturut $\sim$.
Dadi $s = \rho_1(r_1) = \rho_2(r_1)$, lan mula
$\rho_1 = \rho_2$.
\end{proof}

Saiki kasunyatan aljabar mau kita trapaké marang bunder:

\begin{lem}\ollabel{pseudobanachtarski}
Ana partisi $\onesphere$ dadi rong himpunan pisah, $D_{1}$
lan $D_{2}$, saéngga $D_{1}$ bisa dipérang dadi himpunan-himpunan
kaétung cacahé kang bisa diputer kanggo mbentuk salinan
$\onesphere$; semono uga $D_{2}$.
\end{lem}

\begin{proof}
Kanthi $\rotationsgroup_{1}$ lan $\rotationsgroup_{2}$ saka
\olref{disjointgroup}, tetepna:
\begin{align*}
	D_{1} &= \bigcup_{\rho \in \rotationsgroup_1} \funimage{\rho}{C} &
	D_{2} &= \bigcup_{\rho \in \rotationsgroup_2} \funimage{\rho}{C}
\end{align*}
Iki partisi $\onesphere$ miturut \olref{vitalicover}, lan $D_1$
lan $D_2$ pisah miturut \olref{vitalinooverlap}. Miturut konstruksi,
$D_1$ bisa dipérang dadi himpunan-himpunan kaétung cacahé,
$\funimage{\rho}{C}$ kanggo saben $\rho \in \rotationsgroup_1$.
Himpunan-himpunan iki bisa diputer kanggo mbentuk salinan
$\onesphere$, amarga $\onesphere = \bigcup_{\rho \in
\rotationsgroup}\funimage{\rho}{C} = \bigcup_{\rho \in
\rotationsgroup_1}\funimage{(\comp{\rho}{\rho})}{C}$ miturut
\olref{disjointgroup} lan \olref{vitalicover}. Panalaran kang
padha lumaku kanggo $D_2$. \end{proof}\noindent Iki langsung
mènèhi Paradoks Vitali. Kita bisa ngasilaké \emph{rong} salinan
$\onesphere$ saka $\onesphere$ kanthi mbagi dadi pérangan kaétung
cacahé (yaiku $\funimage{\rho}{C}$ kang warna-warna) banjur
mindhah kanthi gerak kaku: saben pérangan $D_1$ cukup diputer,
déné saben pérangan $D_2$ dhisik ditranslasèkaké banjur diputer.

Ayo kita ringkes siasat buktiné. Kita miwiti saka kasunyatan aljabar
grup rotasi ing bidang. Grup iki digunakaké kanggo mbagi
$\onesphere$ dadi kelas-kelas ekuivalensi. Banjur muncul
``paradoks'' nalika Pilihan digunakaké kanggo milih unsur saka
saben kelas.

Siasat kang padha uga digunakaké kanggo mbuktekaké Banach--Tarski.
Béda utamané yaiku kasunyatan aljabar kang dibutuhaké kanggo
Banach--Tarski luwih rumit tinimbang kanggo Paradoks Vitali.
Nanging kasunyatan aljaba mau ora gumantung marang Pilihan.
Kita bakal ngringkes kanthi cekak.

Kanggo Banach--Tarski, kita miwiti kanthi mbuktekaké analog saka
\olref{disjointgroup}: sembarang \emph{grup bebas} bisa dipérang
dadi patang pérangan kang, miturut pangertèn intuitif, bisa
``dipindhah'' kanggo mbalèkaké rong salinan grup wutuh.\footnote{
Kasunyatan yèn cukup \emph{patang} pérangan asalé saka
\cite{Robinson1947}. Bukti anyar ana ing
\citet[Téoréma 5.2]{Wagon2016}. Kita manut
\citet[kaca~3]{Weston2003} nalika nyebut iki minangka
``mindhah'' pérangan-pérangan grup.} Banjur kita nuduhaké yèn
rong rotasi tartamtu ngubengi titik asal ing $\Real^3$ ngasilaké
grup rotasi bebas $F$.\footnote{Delengen
\citet[Téoréma 2.1]{Wagon2016}.} (Nganti kéné durung ana Pilihan.)
Saiki rong titik ing lumahing bal diarani ``padha jinisé'' yèn
siji bisa dipindhah menyang sijiné kanthi rotasi ing~$F$.
Banjur kita \emph{migunakaké Pilihan} kanggo milih pas siji titik
saka saben kelas ekuivalensi titik kang ``padha jinisé''.
Kanthi nrapaké pambagian $F$ marang lumahing bal kaya ing
\olref{pseudobanachtarski}, kita mbagi lumahing mau dadi patang
pérangan kang bisa ``dipindhah'' kanggo ngasilaké rong salinan
lumahing bal. Iki mbuktekaké asil \citep{Hausdorff1914}:

\begin{thm}[Paradoks Hausdorff (ing $\ZFC$)]
Lumahing sembarang bal bisa dipérang dadi sawatara pérangan winates
cacahé, banjur disusun maneh kanthi rotasi lan translasi dadi
rong salinan kang pisah saka lumahing bal mau.
\end{thm}

Isih dibutuhaké rong cara aljabar liya kanggo ngasilaké Téoréma
Banach--Tarski kang jangkep, kang ngrembug ora mung lumahing bal
nanging uga njero bal. Nanging iki mung tambahan ing bukti aljabar.
Mula Weston nulis:
\begin{quote}
  [\ldots] asil babagan grup bebas iku \emph{langkah kunci} ing
  bukti Paradoks Banach--Tarski. Saka sudut pandang iki, Paradoks
  Banach--Tarski ora mung pratelan babagan $\Real^3$, nanging
  luwih-luwih pratelan babagan ruweté grup [translasi lan rotasi
  ing $\Real^3$]. \cite[kaca~16]{Weston2003}
\end{quote}
Tegesé, apa kita bisa mbagi dadi pérangan \emph{winates}
(kaya Banach--Tarski) utawa kudu \emph{tanpa wates nanging kaétung}
(kaya Paradoks Vitali), gumantung marang kasunyatan teori grup
kang béda antarané rong dimensi lan telung dimensi.

Pancen, pengamatan pungkasan iki rada ngrusak guyonan ing pungkasan
\olref[banach]{sec}. Amarga tulisan ``Banach--Tarski'' rong dimensi,
tulisan iku kudu dipérang dadi pérangan kang \emph{tanpa wates}
nanging kaétung yèn arep disusun dadi ``Banach--Tarski
Banach--Tarski''. Supaya guyonané bener, tulisana ing telung
dimensi. Iki kita tinggalaké minangka gladhen kanggo pamaos.

Cathetan pungkasan. Ing \olref[banach]{sec}, kita nyebut yèn
``pérangan'' bal kang dikasilaké ora bisa \emph{diukur}, nanging
kaya ``sebaran tanpa wates'' kang ora bisa digambar. Prakara kang
padha uga ana ing panganggoné Pilihan kanggo ngasilaké
\olref{pseudobanachtarski}. Iki perlu dijlentrehaké.

Kita kudu nggambaraké sethithik dhasar teori ukuran (nanging iki
\emph{mung} gambaran; panjenengan bisa maca buku bab
\emph{ukuran}). Ukuran ing himpunan $X$ mènèhi nilai
$\mu(E) \in \Real$ marang saben $E$ ing sawijining
``aljabar-$\sigma$'' ing $X$. Ing kéné kita nganggo ukuran
nonnegatif kang domainé ditutup tumrap rotasi. Rincian liya
ora wigati, kajaba fungsi $\mu$ kudu aditif kaétung:
ukuran gabungan kaétung saka himpunan-himpunan kang pisah iku
gunggung ukurané siji-siji, yaiku
$\mu(\bigcup_{n < \omega} X_n) = \sum_{n < \omega}\mu(X_n)$
yèn $X_n$ pisah sepasang-sepasang. Himpunan diarani ``ora bisa
diukur'' yèn ora kalebu domain ukuran kang lagi dirembug.
Saiki kita gunakaké $\rotationsgroup$ mau:

\begin{cor}[Vitali]
Anggepen $\mu$ ukuran kang nyukupi $\mu(\onesphere) = 1$,
lan $\mu(X) = \mu(Y)$ yèn $X$ lan $Y$ kongruen. Mula
$\funimage{\rho}{C}$ ora bisa diukur kanggo saben
$\rho \in \rotationsgroup$.
\end{cor}

\begin{proof}
Kanggo bukti kanthi kontradiksi, anggepen salah sawijining bisa
diukur. Jupuken $\mu(\funimage{\sigma}{C}) = r$ kanggo sawijining
$\sigma \in \rotationsgroup$ lan $r \in \Real$. Kanggo sembarang
$\rho \in \rotationsgroup$, himpunan $\funimage{\rho}{C}$ lan
$\funimage{\sigma}{C}$ kongruen, mula
$\mu(\funimage{\rho}{C}) = r$ kanggo saben
$\rho \in \rotationsgroup$. Miturut \olref{vitalicover} lan
\olref{vitalinooverlap}, $\onesphere =
\bigcup_{\rho \in \rotationsgroup}\funimage{\rho}{C}$ iku
gabungan kaétung saka himpunan-himpunan kang pisah. Mula
aditivitas kaétung mènèhi yèn $\mu(\onesphere) = 1$ iku
gunggung ukuran kabèh $\funimage{\rho}{C}$, yaiku
\[
	1 = \mu(\onesphere) = \sum_{\rho \in \rotationsgroup}\mu(\funimage{\rho}{C}) = \sum_{\rho \in \rotationsgroup}r
\]
Nanging yèn $r = 0$, mula $\sum_{\rho \in \rotationsgroup}r = 0$;
yèn $r > 0$, mula $\sum_{\rho \in \rotationsgroup}r = \infty$.
Kaloroné nalisir angka siji ing sisih kiwa.
\end{proof}

\end{document}
